% --- Archon physics formalization source begin ---
% archon:physics
% archon:covers PhyXMiniProblems/problem_phyx_mini_0806.lean
% archon:source-report reports/phyx_mini/problem_phyx_mini_0806.source.json

\chapter{Physics problem phyx\_mini\_0806}
\label{ch:PhyXMiniProblems_problem_phyx_mini_0806}

\paragraph{Problem source.}
\#\# Physical scenario

A sled with a mass of \textbackslash{}( 25.0 \textbackslash{}mathrm\{kg\} \textbackslash{}) rests on a horizontal sheet of essentially frictionless ice. It is attached by a \textbackslash{}( 5.00 \textbackslash{}mathrm\{m\} \textbackslash{}) rope to a post set in the ice. Once given a push, the sled revolves uniformly in a circle around the post as shown in figure

\#\# Question

If the sled makes five complete revolutions every minute, find the force \textbackslash{}( F \textbackslash{}) exerted on it by the rope.

\#\# Answer choices

- A: A: 35.6N

- B: B: 34.3N

- C: C: 32.4N

- D: D: 39.9N

\#\# Figure caption (auxiliary; use the image as primary evidence)

The image depicts a schematic diagram on a flat, reflective surface. There is a blue dashed circular path with arrows indicating a clockwise direction. At the center of the circle, a vertical pole extends upward. An orange sled is placed on the circle's perimeter, connected to the pole by a black line, which appears to denote a radius labeled "R" with an arrow pointing outward from the pole.

\#\# Dataset metadata

Dynamics; Physical Model Grounding Reasoning, Multi-Formula Reasoning

\paragraph{Recorded answer/context.}
B: B: 34.3N

\paragraph{Figure/image path.}
/root/proposal\_for\_physic/science-mango/phyx\_mini\_run/phyx\_data/test\_image/806.png

\paragraph{Formalization target.}
create a compiling Lean file with sorry bodies at `PhyXMiniProblems/problem\_phyx\_mini\_0806.lean`. The Lean declarations must preserve the physical quantities, dimensions or dimensional roles, figure labels, governing-law hypotheses, and final relation expressed by this problem.
Use Mathlib/Physlib names found through LeanExplore where available. If a physics API is missing, introduce faithful local abstractions rather than scalar placeholder aliases.

\begin{theorem}[Physics formalization target]
\label{thm:physics:phyx_mini_0806:target}
\lean{PhyXMiniProblems.ProblemPhyXMini0806.sledRopeForce_is_34_point_3_newtons}
\uses{def:physics:phyx-mini-0806:phyxminiproblems-problemphyxmini0806-forceinnewtons, def:physics:phyx-mini-0806:phyxminiproblems-problemphyxmini0806-sledorbitsetup, def:physics:phyx-mini-0806:phyxminiproblems-problemphyxmini0806-matchessledscenario, def:physics:phyx-mini-0806:phyxminiproblems-problemphyxmini0806-matchesproblemreadouts, def:physics:phyx-mini-0806:phyxminiproblems-problemphyxmini0806-matchesprimaryfigure, def:physics:phyx-mini-0806:phyxminiproblems-problemphyxmini0806-satisfiestautropegeometry, def:physics:phyx-mini-0806:phyxminiproblems-problemphyxmini0806-satisfiesangularrateconversion, def:physics:phyx-mini-0806:phyxminiproblems-problemphyxmini0806-satisfiesuniformcircularkinematics, def:physics:phyx-mini-0806:phyxminiproblems-problemphyxmini0806-satisfieshorizontalnewtoniandynamics, def:physics:phyx-mini-0806:phyxminiproblems-problemphyxmini0806-displayedforceinnewtons, def:physics:phyx-mini-0806:phyxminiproblems-problemphyxmini0806-isclosestdisplayedanswer}
The assigned autoformalize agent should translate this physics problem into Lean declarations in the covered file, with theorem and lemma proof bodies written as `by sorry`.
\end{theorem}
\begin{proof}
This is an autoformalization task, not a proof task. Produce statements that can later be proved by the physics prover without weakening the physical model.
\end{proof}

% --- Archon declaration topology begin ---
\section{Declaration topology}
\begin{definition}[acceleration Dimension]
  \label{def:physics:phyx-mini-0806:phyxminiproblems-problemphyxmini0806-accelerationdimension}
  \lean{PhyXMiniProblems.ProblemPhyXMini0806.accelerationDimension}
  The physical dimension of acceleration, L T⁻²
\end{definition}
\begin{proof}
  This is the declared model definition of acceleration Dimension.
\end{proof}
\begin{definition}[force Dimension]
  \label{def:physics:phyx-mini-0806:phyxminiproblems-problemphyxmini0806-forcedimension}
  \lean{PhyXMiniProblems.ProblemPhyXMini0806.forceDimension}
  \uses{def:physics:phyx-mini-0806:phyxminiproblems-problemphyxmini0806-accelerationdimension}
  The physical dimension of force, M L T⁻²
\end{definition}
\begin{proof}
  This is the declared model definition of force Dimension.
\end{proof}
\begin{definition}[Mass Quantity]
  \label{def:physics:phyx-mini-0806:phyxminiproblems-problemphyxmini0806-massquantity}
  \lean{PhyXMiniProblems.ProblemPhyXMini0806.MassQuantity}
  A nonnegative, unit-independent physical mass
\end{definition}
\begin{proof}
  This is the declared model definition of Mass Quantity.
\end{proof}
\begin{definition}[Length Quantity]
  \label{def:physics:phyx-mini-0806:phyxminiproblems-problemphyxmini0806-lengthquantity}
  \lean{PhyXMiniProblems.ProblemPhyXMini0806.LengthQuantity}
  A nonnegative physical length, used for the rope length and orbit radius
\end{definition}
\begin{proof}
  This is the declared model definition of Length Quantity.
\end{proof}
\begin{definition}[Revolution Rate Quantity]
  \label{def:physics:phyx-mini-0806:phyxminiproblems-problemphyxmini0806-revolutionratequantity}
  \lean{PhyXMiniProblems.ProblemPhyXMini0806.RevolutionRateQuantity}
  A nonnegative revolution rate, with physical dimension inverse time
\end{definition}
\begin{proof}
  This is the declared model definition of Revolution Rate Quantity.
\end{proof}
\begin{definition}[Angular Speed Quantity]
  \label{def:physics:phyx-mini-0806:phyxminiproblems-problemphyxmini0806-angularspeedquantity}
  \lean{PhyXMiniProblems.ProblemPhyXMini0806.AngularSpeedQuantity}
  A nonnegative angular-speed magnitude; radians are dimensionless
\end{definition}
\begin{proof}
  This is the declared model definition of Angular Speed Quantity.
\end{proof}
\begin{definition}[Acceleration Quantity]
  \label{def:physics:phyx-mini-0806:phyxminiproblems-problemphyxmini0806-accelerationquantity}
  \lean{PhyXMiniProblems.ProblemPhyXMini0806.AccelerationQuantity}
  \uses{def:physics:phyx-mini-0806:phyxminiproblems-problemphyxmini0806-accelerationdimension}
  A nonnegative centripetal-acceleration magnitude
\end{definition}
\begin{proof}
  This is the declared model definition of Acceleration Quantity.
\end{proof}
\begin{definition}[Force Quantity]
  \label{def:physics:phyx-mini-0806:phyxminiproblems-problemphyxmini0806-forcequantity}
  \lean{PhyXMiniProblems.ProblemPhyXMini0806.ForceQuantity}
  \uses{def:physics:phyx-mini-0806:phyxminiproblems-problemphyxmini0806-forcedimension}
  A nonnegative force magnitude, whose coherent SI unit is the newton
\end{definition}
\begin{proof}
  This is the declared model definition of Force Quantity.
\end{proof}
\begin{definition}[nonnegative SIReadout]
  \label{def:physics:phyx-mini-0806:phyxminiproblems-problemphyxmini0806-nonnegativesireadout}
  \lean{PhyXMiniProblems.ProblemPhyXMini0806.nonnegativeSIReadout}
  Read any nonnegative dimensionful quantity in coherent SI units
\end{definition}
\begin{proof}
  This is the declared model definition of nonnegative SIReadout.
\end{proof}
\begin{definition}[mass In Kilograms]
  \label{def:physics:phyx-mini-0806:phyxminiproblems-problemphyxmini0806-massinkilograms}
  \lean{PhyXMiniProblems.ProblemPhyXMini0806.massInKilograms}
  \uses{def:physics:phyx-mini-0806:phyxminiproblems-problemphyxmini0806-massquantity, def:physics:phyx-mini-0806:phyxminiproblems-problemphyxmini0806-nonnegativesireadout}
  Kilogram readout of a physical mass
\end{definition}
\begin{proof}
  This is the declared model definition of mass In Kilograms.
\end{proof}
\begin{definition}[length In Meters]
  \label{def:physics:phyx-mini-0806:phyxminiproblems-problemphyxmini0806-lengthinmeters}
  \lean{PhyXMiniProblems.ProblemPhyXMini0806.lengthInMeters}
  \uses{def:physics:phyx-mini-0806:phyxminiproblems-problemphyxmini0806-lengthquantity, def:physics:phyx-mini-0806:phyxminiproblems-problemphyxmini0806-nonnegativesireadout}
  Metre readout of a physical length
\end{definition}
\begin{proof}
  This is the declared model definition of length In Meters.
\end{proof}
\begin{definition}[revolution Rate In Hertz]
  \label{def:physics:phyx-mini-0806:phyxminiproblems-problemphyxmini0806-revolutionrateinhertz}
  \lean{PhyXMiniProblems.ProblemPhyXMini0806.revolutionRateInHertz}
  \uses{def:physics:phyx-mini-0806:phyxminiproblems-problemphyxmini0806-revolutionratequantity, def:physics:phyx-mini-0806:phyxminiproblems-problemphyxmini0806-nonnegativesireadout}
  Cycle-per-second readout of a physical revolution rate
\end{definition}
\begin{proof}
  This is the declared model definition of revolution Rate In Hertz.
\end{proof}
\begin{definition}[angular Speed In Radians Per Second]
  \label{def:physics:phyx-mini-0806:phyxminiproblems-problemphyxmini0806-angularspeedinradianspersecond}
  \lean{PhyXMiniProblems.ProblemPhyXMini0806.angularSpeedInRadiansPerSecond}
  \uses{def:physics:phyx-mini-0806:phyxminiproblems-problemphyxmini0806-angularspeedquantity, def:physics:phyx-mini-0806:phyxminiproblems-problemphyxmini0806-nonnegativesireadout}
  Radian-per-second readout of a physical angular-speed magnitude
\end{definition}
\begin{proof}
  This is the declared model definition of angular Speed In Radians Per Second.
\end{proof}
\begin{definition}[acceleration In Meters Per Second Squared]
  \label{def:physics:phyx-mini-0806:phyxminiproblems-problemphyxmini0806-accelerationinmeterspersecondsquared}
  \lean{PhyXMiniProblems.ProblemPhyXMini0806.accelerationInMetersPerSecondSquared}
  \uses{def:physics:phyx-mini-0806:phyxminiproblems-problemphyxmini0806-accelerationquantity, def:physics:phyx-mini-0806:phyxminiproblems-problemphyxmini0806-nonnegativesireadout}
  Metre-per-second-squared readout of an acceleration magnitude
\end{definition}
\begin{proof}
  This is the declared model definition of acceleration In Meters Per Second Squared.
\end{proof}
\begin{definition}[force In Newtons]
  \label{def:physics:phyx-mini-0806:phyxminiproblems-problemphyxmini0806-forceinnewtons}
  \lean{PhyXMiniProblems.ProblemPhyXMini0806.forceInNewtons}
  \uses{def:physics:phyx-mini-0806:phyxminiproblems-problemphyxmini0806-forcequantity, def:physics:phyx-mini-0806:phyxminiproblems-problemphyxmini0806-nonnegativesireadout}
  Newton readout of a force magnitude
\end{definition}
\begin{proof}
  This is the declared model definition of force In Newtons.
\end{proof}
\begin{definition}[Trajectory Shape]
  \label{def:physics:phyx-mini-0806:phyxminiproblems-problemphyxmini0806-trajectoryshape}
  \lean{PhyXMiniProblems.ProblemPhyXMini0806.TrajectoryShape}
  Shape of the sled's trajectory about the post
\end{definition}
\begin{proof}
  This is the declared model definition of Trajectory Shape.
\end{proof}
\begin{definition}[Surface Orientation]
  \label{def:physics:phyx-mini-0806:phyxminiproblems-problemphyxmini0806-surfaceorientation}
  \lean{PhyXMiniProblems.ProblemPhyXMini0806.SurfaceOrientation}
  Orientation of the ice sheet
\end{definition}
\begin{proof}
  This is the declared model definition of Surface Orientation.
\end{proof}
\begin{definition}[Friction Regime]
  \label{def:physics:phyx-mini-0806:phyxminiproblems-problemphyxmini0806-frictionregime}
  \lean{PhyXMiniProblems.ProblemPhyXMini0806.FrictionRegime}
  Tangential resistance regime of the ice
\end{definition}
\begin{proof}
  This is the declared model definition of Friction Regime.
\end{proof}
\begin{definition}[Radial Direction]
  \label{def:physics:phyx-mini-0806:phyxminiproblems-problemphyxmini0806-radialdirection}
  \lean{PhyXMiniProblems.ProblemPhyXMini0806.RadialDirection}
  Radial direction of the force exerted by the rope on the sled
\end{definition}
\begin{proof}
  This is the declared model definition of Radial Direction.
\end{proof}
\begin{definition}[Rotation Sense]
  \label{def:physics:phyx-mini-0806:phyxminiproblems-problemphyxmini0806-rotationsense}
  \lean{PhyXMiniProblems.ProblemPhyXMini0806.RotationSense}
  Sense of the motion indicated by the blue arrowheads in the bitmap
\end{definition}
\begin{proof}
  This is the declared model definition of Rotation Sense.
\end{proof}
\begin{definition}[Horizontal Force Source]
  \label{def:physics:phyx-mini-0806:phyxminiproblems-problemphyxmini0806-horizontalforcesource}
  \lean{PhyXMiniProblems.ProblemPhyXMini0806.HorizontalForceSource}
  Possible sources of the net horizontal force on the sled
\end{definition}
\begin{proof}
  This is the declared model definition of Horizontal Force Source.
\end{proof}
\begin{definition}[Sled Orbit Figure]
  \label{def:physics:phyx-mini-0806:phyxminiproblems-problemphyxmini0806-sledorbitfigure}
  \lean{PhyXMiniProblems.ProblemPhyXMini0806.SledOrbitFigure}
  \uses{def:physics:phyx-mini-0806:phyxminiproblems-problemphyxmini0806-rotationsense}
  Literal and qualitative evidence transcribed from primary image 806.png
\end{definition}
\begin{proof}
  This is the declared model definition of Sled Orbit Figure.
\end{proof}
\begin{definition}[Sled Orbit Setup]
  \label{def:physics:phyx-mini-0806:phyxminiproblems-problemphyxmini0806-sledorbitsetup}
  \lean{PhyXMiniProblems.ProblemPhyXMini0806.SledOrbitSetup}
  \uses{def:physics:phyx-mini-0806:phyxminiproblems-problemphyxmini0806-massquantity, def:physics:phyx-mini-0806:phyxminiproblems-problemphyxmini0806-lengthquantity, def:physics:phyx-mini-0806:phyxminiproblems-problemphyxmini0806-revolutionratequantity, def:physics:phyx-mini-0806:phyxminiproblems-problemphyxmini0806-angularspeedquantity, def:physics:phyx-mini-0806:phyxminiproblems-problemphyxmini0806-accelerationquantity, def:physics:phyx-mini-0806:phyxminiproblems-problemphyxmini0806-forcequantity, def:physics:phyx-mini-0806:phyxminiproblems-problemphyxmini0806-trajectoryshape, def:physics:phyx-mini-0806:phyxminiproblems-problemphyxmini0806-surfaceorientation, def:physics:phyx-mini-0806:phyxminiproblems-problemphyxmini0806-frictionregime, def:physics:phyx-mini-0806:phyxminiproblems-problemphyxmini0806-radialdirection, def:physics:phyx-mini-0806:phyxminiproblems-problemphyxmini0806-horizontalforcesource, def:physics:phyx-mini-0806:phyxminiproblems-problemphyxmini0806-sledorbitfigure}
  Independent physical quantities of the experiment. In particular, ropeForce is neither defined from nor initialized with the displayed answer
\end{definition}
\begin{proof}
  This is the declared model definition of Sled Orbit Setup.
\end{proof}
\begin{definition}[Matches Sled Scenario]
  \label{def:physics:phyx-mini-0806:phyxminiproblems-problemphyxmini0806-matchessledscenario}
  \lean{PhyXMiniProblems.ProblemPhyXMini0806.MatchesSledScenario}
  \uses{def:physics:phyx-mini-0806:phyxminiproblems-problemphyxmini0806-sledorbitsetup}
  Qualitative conditions stated in the physical scenario
\end{definition}
\begin{proof}
  This is the declared model definition of Matches Sled Scenario.
\end{proof}
\begin{definition}[Matches Problem Readouts]
  \label{def:physics:phyx-mini-0806:phyxminiproblems-problemphyxmini0806-matchesproblemreadouts}
  \lean{PhyXMiniProblems.ProblemPhyXMini0806.MatchesProblemReadouts}
  \uses{def:physics:phyx-mini-0806:phyxminiproblems-problemphyxmini0806-massinkilograms, def:physics:phyx-mini-0806:phyxminiproblems-problemphyxmini0806-lengthinmeters, def:physics:phyx-mini-0806:phyxminiproblems-problemphyxmini0806-revolutionrateinhertz, def:physics:phyx-mini-0806:phyxminiproblems-problemphyxmini0806-sledorbitsetup}
  Numerical data from the prose. Five revolutions per minute is recorded as 5 / 60 cycles per coherent-SI second. No force value occurs here
\end{definition}
\begin{proof}
  This is the declared model definition of Matches Problem Readouts.
\end{proof}
\begin{definition}[Matches Primary Figure]
  \label{def:physics:phyx-mini-0806:phyxminiproblems-problemphyxmini0806-matchesprimaryfigure}
  \lean{PhyXMiniProblems.ProblemPhyXMini0806.MatchesPrimaryFigure}
  \uses{def:physics:phyx-mini-0806:phyxminiproblems-problemphyxmini0806-sledorbitsetup}
  Primary-image evidence. The bitmap's lower arrow points to the right and its upper arrow points to the left, so its displayed sense is counterclockwise; this follows the required primary evidence rather than the auxiliary caption's incompatible word "clockwise"
\end{definition}
\begin{proof}
  This is the declared model definition of Matches Primary Figure.
\end{proof}
\begin{definition}[Satisfies Taut Rope Geometry]
  \label{def:physics:phyx-mini-0806:phyxminiproblems-problemphyxmini0806-satisfiestautropegeometry}
  \lean{PhyXMiniProblems.ProblemPhyXMini0806.SatisfiesTautRopeGeometry}
  \uses{def:physics:phyx-mini-0806:phyxminiproblems-problemphyxmini0806-sledorbitsetup}
  A taut radial rope identifies its physical length with orbit radius R
\end{definition}
\begin{proof}
  This is the declared model definition of Satisfies Taut Rope Geometry.
\end{proof}
\begin{definition}[Satisfies Angular Rate Conversion]
  \label{def:physics:phyx-mini-0806:phyxminiproblems-problemphyxmini0806-satisfiesangularrateconversion}
  \lean{PhyXMiniProblems.ProblemPhyXMini0806.SatisfiesAngularRateConversion}
  \uses{def:physics:phyx-mini-0806:phyxminiproblems-problemphyxmini0806-revolutionrateinhertz, def:physics:phyx-mini-0806:phyxminiproblems-problemphyxmini0806-angularspeedinradianspersecond, def:physics:phyx-mini-0806:phyxminiproblems-problemphyxmini0806-sledorbitsetup}
  One complete revolution is 2π radians
\end{definition}
\begin{proof}
  This is the declared model definition of Satisfies Angular Rate Conversion.
\end{proof}
\begin{definition}[Satisfies Uniform Circular Kinematics]
  \label{def:physics:phyx-mini-0806:phyxminiproblems-problemphyxmini0806-satisfiesuniformcircularkinematics}
  \lean{PhyXMiniProblems.ProblemPhyXMini0806.SatisfiesUniformCircularKinematics}
  \uses{def:physics:phyx-mini-0806:phyxminiproblems-problemphyxmini0806-lengthinmeters, def:physics:phyx-mini-0806:phyxminiproblems-problemphyxmini0806-angularspeedinradianspersecond, def:physics:phyx-mini-0806:phyxminiproblems-problemphyxmini0806-accelerationinmeterspersecondsquared, def:physics:phyx-mini-0806:phyxminiproblems-problemphyxmini0806-sledorbitsetup}
  Uniform circular kinematics gives centripetal acceleration a = ω² R
\end{definition}
\begin{proof}
  This is the declared model definition of Satisfies Uniform Circular Kinematics.
\end{proof}
\begin{definition}[Satisfies Horizontal Newtonian Dynamics]
  \label{def:physics:phyx-mini-0806:phyxminiproblems-problemphyxmini0806-satisfieshorizontalnewtoniandynamics}
  \lean{PhyXMiniProblems.ProblemPhyXMini0806.SatisfiesHorizontalNewtonianDynamics}
  \uses{def:physics:phyx-mini-0806:phyxminiproblems-problemphyxmini0806-massinkilograms, def:physics:phyx-mini-0806:phyxminiproblems-problemphyxmini0806-accelerationinmeterspersecondsquared, def:physics:phyx-mini-0806:phyxminiproblems-problemphyxmini0806-forceinnewtons, def:physics:phyx-mini-0806:phyxminiproblems-problemphyxmini0806-sledorbitsetup}
  On frictionless horizontal ice, rope tension is the only horizontal force and Newton's second law in the inward radial direction is F = m a. This is a general governing relation and contains no requested numerical force
\end{definition}
\begin{proof}
  This is the declared model definition of Satisfies Horizontal Newtonian Dynamics.
\end{proof}
\begin{definition}[Answer Choice]
  \label{def:physics:phyx-mini-0806:phyxminiproblems-problemphyxmini0806-answerchoice}
  \lean{PhyXMiniProblems.ProblemPhyXMini0806.AnswerChoice}
  Labels of the four force values displayed in the question
\end{definition}
\begin{proof}
  This is the declared model definition of Answer Choice.
\end{proof}
\begin{definition}[displayed Force In Newtons]
  \label{def:physics:phyx-mini-0806:phyxminiproblems-problemphyxmini0806-displayedforceinnewtons}
  \lean{PhyXMiniProblems.ProblemPhyXMini0806.displayedForceInNewtons}
  \uses{def:physics:phyx-mini-0806:phyxminiproblems-problemphyxmini0806-answerchoice}
  Displayed answer values, read in newtons
\end{definition}
\begin{proof}
  This is the declared model definition of displayed Force In Newtons.
\end{proof}
\begin{definition}[Is Closest Displayed Answer]
  \label{def:physics:phyx-mini-0806:phyxminiproblems-problemphyxmini0806-isclosestdisplayedanswer}
  \lean{PhyXMiniProblems.ProblemPhyXMini0806.IsClosestDisplayedAnswer}
  \uses{def:physics:phyx-mini-0806:phyxminiproblems-problemphyxmini0806-answerchoice, def:physics:phyx-mini-0806:phyxminiproblems-problemphyxmini0806-displayedforceinnewtons}
  A displayed choice is closest to a computed force magnitude
\end{definition}
\begin{proof}
  This is the declared model definition of Is Closest Displayed Answer.
\end{proof}
% --- Archon declaration topology end ---
% --- Archon physics formalization source end ---
